\documentclass[hyperref={bookmarks=false},aspectratio=169]{beamer}
\usepackage{hyperref,ragged2e,amsmath,amssymb,url,bm}
\usepackage{graphicx,booktabs,array}
\usepackage{tikz}
\usetikzlibrary{arrows.meta,positioning,calc,shapes.geometric,decorations.pathreplacing,matrix}
\graphicspath{{crops/}}
\usetheme{Goettingen}
% --- Paleta "Israel": celeste y azul del tejelet (Nm 15:38), oro del santuario, fondo crema ---
% (los nombres de color se conservan por compatibilidad con los cuerpos de las láminas: «ladrillo» = azul tejelet,
%  «cafe» = azul real, «cafeosc» = azul noche, «piedra» = celeste claro, «sepia» = azul pizarra)
\definecolor{cafe}{RGB}{26,82,150}
\definecolor{cafeosc}{RGB}{11,33,74}
\definecolor{tinta}{RGB}{20,32,56}
\definecolor{oro}{RGB}{206,162,52}
\definecolor{ladrillo}{RGB}{20,68,140}
\definecolor{piedra}{RGB}{234,243,251}
\definecolor{sepia}{RGB}{72,102,140}
\definecolor{celeste}{RGB}{96,166,222}
\setbeamercolor{structure}{fg=ladrillo}
\setbeamertemplate{sidebar canvas right}[vertical shading][top=celeste!75!ladrillo,middle=ladrillo,bottom=cafeosc]
\setbeamercolor{title in sidebar}{fg=oro!80!white}
\setbeamercolor{author in sidebar}{fg=piedra}
\setbeamercolor{section in sidebar}{fg=white}
\setbeamercolor{section in sidebar shaded}{fg=celeste!45!white}
\setbeamercolor{subsection in sidebar}{fg=white}
\setbeamercolor{subsection in sidebar shaded}{fg=celeste!45!white}
\let\raggedright\justifying
\definecolor{crema}{RGB}{250,245,230}
\setbeamertemplate{background canvas}[default]
\setbeamercolor{background canvas}{bg=crema}
\setbeamercolor{normal text}{fg=tinta}
\setbeamercolor{frametitle}{fg=ladrillo}
\setbeamerfont{frametitle}{series=\bfseries}
\setbeamertemplate{section in toc shaded}{\leavevmode{\color{sepia!85!black}\inserttocsection}\par}
\setbeamercolor{title}{fg=ladrillo}
\setbeamercolor{subtitle}{fg=cafe}
\setbeamercolor{item projected}{bg=cafe,fg=piedra}
\setbeamertemplate{itemize item}{\small\textcolor{oro!85!ladrillo}{$\blacktriangleright$}}
\setbeamertemplate{navigation symbols}{}
% --- Español + hebreo + griego (XeLaTeX) con fuentes empaquetadas en fonts/ ---
\usepackage[bidi=default]{babel}
\babelprovide[import,main]{spanish}
\babelprovide[import]{hebrew}
\babelprovide[import]{ancientgreek}
\babelfont[hebrew]{sf}[Path=fonts/, BoldFont=NotoHebrewSans-700.ttf]{NotoHebrewSans-400.ttf}
\babelfont[hebrew]{rm}[Path=fonts/, BoldFont=NotoHebrewSans-700.ttf]{NotoHebrewSans-400.ttf}
\babelfont[ancientgreek]{sf}[Path=fonts/]{GFSPorson.otf}
\babelfont[ancientgreek]{rm}[Path=fonts/]{GFSPorson.otf}
\newcommand{\heb}[1]{\mbox{\foreignlanguage{hebrew}{#1}}}
\newcommand{\grc}[1]{\foreignlanguage{ancientgreek}{#1}}
% --- Citas, imágenes, referencias ---
\usepackage[most]{tcolorbox}
\newtcolorbox{cita}[1][]{enhanced, colback=piedra!92!white, colframe=oro!90!ladrillo,
  boxrule=0pt, leftrule=2.5pt, arc=0pt, left=6pt, right=6pt, top=4pt, bottom=4pt,
  before skip=4pt, after skip=6pt, fontupper=\small\itshape, #1}
\newcommand{\fuente}[1]{\par\hfill{\upshape\bfseries\footnotesize\textcolor{ladrillo}{#1}}}
\newcommand{\imagen}[3][0.80]{%
  \centering\includegraphics[height=#1\textheight,width=\linewidth,keepaspectratio]{#2}\par\vspace{2pt}
  {\fontsize{5}{6}\selectfont\color{sepia}{\raggedleft\textit{Nota.} #3\par}}}
\makeatletter
\newcommand{\fotolado}[2]{%
  \begin{tikzpicture}[remember picture,overlay]
    \node[anchor=north east,inner sep=0pt] (foto) at ([xshift=-\beamer@sidebarwidth]current page.north east)
      {\includegraphics[height=\paperheight]{#1}};
    \node[anchor=south east,inner sep=2pt,text=piedra,fill=cafeosc,fill opacity=0.6,text opacity=1,
      font=\fontsize{4.5}{5.5}\selectfont,align=right,text width=4.2cm] at ([xshift=-3pt,yshift=3pt]foto.south east)
      {\textit{Nota.} #2};
  \end{tikzpicture}}
\makeatother
\newcommand{\refapa}[1]{\par\hangindent=1.2em\hangafter=1\noindent #1\par\vspace{3pt}}
\usepackage{xurl}\urlstyle{same}
\newcommand{\destacar}[1]{\textbf{\textcolor{ladrillo}{#1}}}
\usetikzlibrary{shadows}
\tikzset{caja/.style={draw=cafe!80!black,line width=0.7pt,top color=piedra!55!white,bottom color=piedra!92!oro,
  rounded corners=4pt,inner sep=5pt,font=\scriptsize,
  drop shadow={shadow xshift=0.7pt,shadow yshift=-0.8pt,opacity=0.25,fill=cafeosc}},
  >={Stealth[length=2.4mm,width=1.8mm]}, every picture/.style={line cap=round,line join=round}}
\newcommand{\precision}[1]{{\footnotesize\color{sepia}\textbf{Precisión:} #1}}
\newcommand{\pt}[2]{({#1-30},{#2-29})}
